\chapter{Краткий справочник API}
\label{app:api}

Справочник перечисляет основные публичные имена из пакета \code{ciao}. Точные
ограничения сигнатур и ошибок определяются актуальной технической спецификацией.

\section{Декларация модели}

\small
\begin{tabularx}{\textwidth}{@{}p{0.36\textwidth}X@{}}
\toprule
Имя & Назначение \\
\midrule
\code{Automaton[Context]} & Базовый класс декларативного автомата. \\
\code{ControllerAutomaton} & Роль автомата с доменной управляющей логикой. \\
\code{ViewAutomaton} & Роль автомата-посредника между Controller и UI. \\
\code{AutomatonContext} & База для явно передаваемого контекста. \\
\code{State(initial, final)} & Дескриптор состояния. \\
\code{AutomatonObject(...)} & Именованное включение автомата. \\
\code{Program(...)} & Полная декларация исполняемой программы. \\
\code{Connection(required, provided)} & Маршрут совместимых операций. \\
\bottomrule
\end{tabularx}

\section{Декораторы и переходы}

\begin{tabularx}{\textwidth}{@{}p{0.36\textwidth}X@{}}
\toprule
Имя & Назначение \\
\midrule
\code{@provided_command} & Входная асинхронная операция владельца. \\
\code{@required_command} & Исходящая асинхронная зависимость. \\
\code{@provided_query} & Синхронная реализация ответа. \\
\code{@required_query} & Синхронная зависимость вызывающего объекта. \\
\code{@guard} & Булево условие выбора перехода. \\
\code{@action} & Эффект перехода; может объявлять query/emit capabilities. \\
\code{transition(...)} & Связывает source, event, guard/actions и target. \\
\bottomrule
\end{tabularx}

\section{Внедряемые capabilities}

\begin{tabularx}{\textwidth}{@{}p{0.36\textwidth}X@{}}
\toprule
Тип & Назначение \\
\midrule
\code{Emit} & Вызов одной declared required-команды. \\
\code{Emits} & Индексируемый набор declared required-команд. \\
\code{Query} & Вызов одного declared required-query. \\
\code{Queries} & Индексируемый набор declared required-queries. \\
\bottomrule
\end{tabularx}

\section{Исполнение}

\begin{tabularx}{\textwidth}{@{}p{0.36\textwidth}X@{}}
\toprule
Имя & Назначение \\
\midrule
\code{compile_program(program)} & Валидирует и компилирует декларации. \\
\code{Runtime(...)} & Создаёт изолированный исполнитель. \\
\code{Runtime.send(...)} & Кладёт external provided-команду в FIFO. \\
\code{Runtime.step()} & Обрабатывает не более одного сообщения. \\
\code{Runtime.run(max_steps=...)} & Обрабатывает очередь до остановки. \\
\code{Runtime.snapshot(object)} & Возвращает текущий снимок объекта. \\
\code{Runtime.trace} & Неизменяемая последовательность событий исполнения. \\
\bottomrule
\end{tabularx}

\section{UI}

\begin{tabularx}{\textwidth}{@{}p{0.36\textwidth}X@{}}
\toprule
Имя & Назначение \\
\midrule
\code{TkConnection} & Адаптер Button/endpoints/handlers для Tkinter. \\
\code{FactoryPayloadSource} & Создаёт payload по данным Window в момент события. \\
\code{ConstantPayloadSource} & Передаёт заранее заданный payload. \\
\code{Subscription.dispose()} & Снимает регистрацию или освобождает адаптер. \\
\bottomrule
\end{tabularx}

\section{Минимальная памятка}

\begin{enumerate}
  \item Объявите Context, Automaton, State, operations и transitions.
  \item Создайте каждый \code{AutomatonObject} с уникальным именем.
  \item Соедините required endpoints с provided endpoints.
  \item Создайте Runtime; для ранней диагностики можно отдельно вызвать
    \code{compile_program}.
  \item Отправляйте только external provided-команды полным endpoint.
  \item Анализируйте \code{StepResult}, \code{RunResult} и \code{trace}.
\end{enumerate}

